\subsection{WEIGHTS OF A LINEAR REPRESENTATION}

In this number k denotes one of the fields R or C.

Let V be a finite dimensional vector space over k, and \(\rho: G \to \mathbf{GL}(V)\) a continuous (hence real-analytic, Chap. III, §8, no. 1, Th. 1) representation of the connected compact Lie group G on V. Define a complex vector space \(\tilde{V}\) and a continuous representation \(\tilde{\rho}: G \to \mathbf{GL}(\tilde{V})\) as follows: if k = C, set \(\tilde{V} = V\) , \(\tilde{\rho} = \rho\) ; if k = R, set \(\tilde{V} = V_{(C)}\) and \(\tilde{\rho}\) to be the composite of \(\rho\) with the canonical homomorphism \(\mathbf{GL}(V) \to \mathbf{GL}(\tilde{V})\) .

For all \(\lambda \in \mathrm{X}(\mathrm{G})\) , denote by \(\tilde{\mathrm{V}}_{\lambda}(\mathrm{G})\) the vector subspace of \(\tilde{\mathrm{V}}\) consisting of the \(v \in \tilde{\mathrm{V}}\) such that \(\tilde{\rho}(g)v = g^{\lambda}v\) for all \(g \in \mathrm{G}\) (cf.~Chap. VII, §1, no. 1). By loc. cit., Prop. 3, the sum of the \(\tilde{\mathrm{V}}_{\lambda}(\mathrm{G})\) (for \(\lambda\) belonging to \(\mathrm{X}(\mathrm{G})\) ) is direct. Moreover:

Lemma 1. If G is commutative, \(\tilde{V}\) is the direct sum of the \(\tilde{\mathrm{V}}_{\lambda}(\mathrm{G})\) for \(\lambda \in \mathrm{X}(\mathrm{G})\) .

Since \(\rho\) is semi-simple (§1, no. 1), it suffices to prove the lemma in the case in which \(\rho\) is simple. In that case, the commutant \(Z\) of \(\rho(G)\) in \(\mathrm{End}(\tilde{V})\) reduces to homotheties (Algebra, Chap. VIII, §3, no. 2, Th. 1); thus, the image of the homomorphism \(\tilde{\rho}\) is contained in the subgroup \(\mathbf{C}^*.1_V\) of \(\mathrm{GL}(\tilde{V})\) , and there exists \(\lambda \in X(G)\) such that \(\tilde{V} = \tilde{V}_{\lambda}(G)\) .

DEFINITION 1. The weights of the representation \(\rho\) of G, relative to a maximal torus T of G, are the elements \(\lambda\) of X(T) such that \(\tilde{\mathrm{V}}_{\lambda}(\mathrm{T}) \neq 0\) .

Denote by \(\mathrm{P}(\rho,\mathrm{T})\) , or by \(\mathrm{P}(\rho)\) if there is no possibility of confusion over the choice of T, the set of weights of \(\rho\) relative to T. By Lemma 1,

\[
\tilde {\mathrm{V}} = \bigoplus_ {\lambda \in \mathrm{P} (\rho , \mathrm{T})} \tilde {\mathrm{V}} _ {\lambda} (\mathrm{T}).\tag{7}
\]

Let \(T'\) be another maximal torus of G and g an element of G such that \((\operatorname{Int} g)T = T'\) (§2, no. 2, Th. 2). For all \(\lambda \in X(T)\) ,

\[
\tilde {\rho} (g) (\tilde {\mathrm{V}} _ {\lambda} (\mathrm{T})) = \tilde {\mathrm{V}} _ {\lambda^ {\prime}} (\mathrm{T} ^ {\prime}), \quad \text { where } \lambda^ {\prime} = \mathrm{X} (\operatorname{Int} g ^ {- 1}) (\lambda).\tag{8}
\]

Consequently,

\[
\mathrm{X(Int} g) (\mathrm{P} (\rho , \mathrm{T} ^ {\prime})) = \mathrm{P} (\rho , \mathrm{T}).\tag{9}
\]

The Weyl group \(\mathrm{W} = \mathrm{W}_{\mathrm{G}}(\mathrm{T})\) operates on the left on the \(\mathbf{Z}\) -module \(\mathrm{X(T)}\) by \(w\mapsto \mathrm{X}(w^{-1})\) ; thus, for \(t\in \mathrm{T},\lambda \in \mathrm{X(T)},w\in \mathrm{W}\) , we have \(t^{w\lambda} = (w^{-1}(t))^{\lambda}\) .

PROPOSITION 5. The set \(\mathrm{P}(\rho, \mathrm{T})\) is stable under the operation of the Weyl group W. Let \(n \in \mathrm{N}_{\mathrm{G}}(\mathrm{T})\) , and let \(w\) be its class in W; for \(\lambda \in \mathrm{X}(\mathrm{T})\) , we have \(\rho(n)(\tilde{\mathrm{V}}_{\lambda}(\mathrm{T})) = \tilde{\mathrm{V}}_{w\lambda}(\mathrm{T})\) and \(\dim \tilde{\mathrm{V}}_{w\lambda}(\mathrm{T}) = \dim \tilde{\mathrm{V}}_{\lambda}(\mathrm{T})\) .

Formula (9), with \(T' = T\) , g = n, implies that \(\mathrm{P}(\rho, \mathrm{T})\) is stable under w; further, \(\tilde{\rho}(n)\) induces an isomorphism from \(\tilde{\mathrm{V}}_{\lambda}(\mathrm{T})\) to \(\tilde{\mathrm{V}}_{w\lambda}(\mathrm{T})\) (formula (8)), hence the proposition.

PROPOSITION 6. The homomorphism \(\rho : \mathbf{G} \to \mathbf{GL}(\mathbf{V})\) is injective if and only if \(\mathrm{P}(\rho, \mathrm{T})\) generates the \(\mathbf{Z}\) -module \(\mathrm{X}(\mathrm{T})\) .

The homomorphism \(\rho\) is injective if and only if its restriction to T is injective (§2, no. 6, Prop. 9). Further, since the canonical homomorphism \(\mathbf{GL}(\mathrm{V}) \to \mathbf{GL}(\tilde{\mathrm{V}})\) is injective, we can replace \(\rho\) by \(\tilde{\rho}\) . It then follows from (7) that the kernel of the restriction of \(\rho\) to T is the intersection of the kernels of the elements of \(\mathrm{P}(\rho, \mathrm{T})\) . Thus, the conclusion follows from Prop. 2 of no. 1.

The linear representation \(\mathrm{L}(\rho)\) of t in \(\mathfrak{gl}(\tilde{\mathrm{V}})\) extends to a homomorphism of C-Lie algebras

\[
\tilde {\mathrm{L}} (\rho): \mathfrak {t} _ {\mathbf {C}} \to \mathfrak {g l} (\tilde {\mathrm{V}}).
\]

Moreover, recall (no. 1) that we have associated to every element \(\lambda\) of X(T) a linear form \(\delta (\lambda)\) on \(\mathbf{t}_{\mathbf{C}}\) such that

\[
(\exp_ {\mathrm{T}} x) ^ {\lambda} = e ^ {\delta (\lambda) (x)}, \qquad x \in \mathfrak {t}.\tag{10}
\]

Recall finally (Chap. VII, §1, no. 1) that, for any map \(\mu : \mathfrak{t}_{\mathbf{C}} \to \mathbf{C}\) , we denote by \(\tilde{\mathrm{V}}_{\mu}(\mathfrak{t}_{\mathbf{C}})\) the vector subspace of \(\tilde{\mathrm{V}}\) consisting of the \(v\) such that \((\tilde{\mathrm{L}}(\rho)(u))(v) = \mu(u).v\) for all \(u \in \mathfrak{t}_{\mathbf{C}}\) .

We now deduce from (7) and loc. cit., Prop. 3:

PROPOSITION 7. a) For all \(\lambda \in \mathrm{X(T)}\) , we have \(\tilde{\mathrm{V}}_{\lambda}(\mathrm{T}) = \tilde{\mathrm{V}}_{\delta (\lambda)}(\mathbf{t}_{\mathbf{C}})\) .

\begin{enumerate}
\def\labelenumi{\alph{enumi})}
\setcounter{enumi}{1}
\tightlist
\item
  The map \(\delta : \mathrm{X(T)} \to \mathrm{Hom}_{\mathbf{C}}(\mathfrak{t}_{\mathbf{C}}, \mathbf{C})\) induces a bijection from \(\mathrm{P}(\rho, \mathrm{T})\) to the set of weights of \(\mathfrak{t}_{\mathbf{C}}\) on \(\tilde{\mathrm{V}}\) .
\end{enumerate}

Note first that, if W operates on \(t_{C}\) by associating to any element w of W the endomorphism \(\mathrm{L}(w)_{(\mathbf{C})}\) of \(t_{C}\) , the map \(\delta\) is compatible with the operation of W on X(T) and \(\mathrm{Hom}_{\mathbf{C}}(t_{\mathbf{C}}, \mathbf{C})\) .

Assume now that \(k = \mathbf{R}\) . Denote by \(\sigma\) the conjugation of \(\tilde{\mathrm{V}}\) relative to \(\mathrm{V}\) , defined by \(\sigma(x + iy) = x - iy\) for \(x, y\) in \(\mathrm{V}\) ; for every complex vector subspace \(\mathrm{E}\) of \(\tilde{\mathrm{V}}\) , the smallest subspace of \(\tilde{\mathrm{V}}\) rational over \(\mathbf{R}\) and containing \(\mathrm{E}\) is \(\mathrm{E} + \sigma(\mathrm{E})\) . In particular, for all \(\lambda \in \mathrm{X}(\mathrm{T})\) , there exists a real vector subspace \(\mathrm{V}(\lambda)\) of \(\mathrm{V}\) such that the subspace \(\mathrm{V}(\lambda)_{(\mathbf{C})}\) of \(\tilde{\mathrm{V}}\) is \(\tilde{\mathrm{V}}_{\lambda}(\mathrm{T}) + \tilde{\mathrm{V}}_{-\lambda}(\mathrm{T})\) (note that \(\sigma(\tilde{\mathrm{V}}_{\lambda}(\mathrm{T})) = \tilde{\mathrm{V}}_{-\lambda}(\mathrm{T})\) ). We have \(\mathrm{V}(\lambda) = \mathrm{V}(-\lambda)\) , and the \(\mathrm{V}(\lambda)\) are the isotypical components of the representation of \(\mathrm{T}\) on \(\mathrm{V}\) induced by \(\rho\) .

